\enlargethispage{2pt}
\Article{30}

Bei einem zusammengesetzten Modul ist manchmal die folgende Methode vorzuziehen.

Der Modul sei $=mn$ und die gegebene Kongruenz $ax\equiv b$. Man löse diese zunächst nach dem Modul $m$ und nehme an, dass sie für $x\equiv v \ (\mathrm{mod.}\ \frac{m}{\delta})$ erfüllt ist, wobei $\delta$ den größten gemeinsamen Teiler von $m,a$ bezeichnet. Jeder Wert von $x$, der die Kongruenz $ax\equiv b$ nach dem Modul $mn$ erfüllt, erfüllt sie offenbar auch nach dem Modul $m$ und hat daher die Form $v+\frac{m}{\delta}x'$, wobei $x'$ eine unbestimmte Zahl bezeichnet. Umgekehrt erfüllen aber nicht alle Zahlen der Form $v+\frac{m}{\delta}x'$ die Kongruenz nach dem Modul $mn$. Wie $x'$ zu bestimmen ist, damit $v+\frac{m}{\delta}x'$ eine Wurzel der Kongruenz $ax\equiv b \ (\mathrm{mod.}\ mn)$ wird, ergibt sich aus der Lösung der Kongruenz $\frac{am}{\delta}x'+av\equiv b\ (\mathrm{mod.}\ mn)$, die mit $\frac{a}{\delta}x'\equiv\frac{b-av}{m}\ (\mathrm{mod.}\ n)$ gleichwertig ist. Somit lässt sich die Lösung jeder Kongruenz ersten Grades nach dem Modul $mn$ auf die Lösung zweier Kongruenzen nach den Moduln $m$ und $n$ zurückführen. Ist $n$ wiederum ein Produkt zweier Faktoren, so hängt die Lösung der Kongruenz nach dem Modul $n$ ebenso von der Lösung zweier Kongruenzen ab, deren Moduln diese Faktoren sind. Allgemein hängt die Lösung einer Kongruenz nach einem beliebigen zusammengesetzten Modul von der Lösung anderer Kongruenzen ab, deren Moduln Faktoren dieser Zahl sind. Diese können, wenn es zweckmäßig erscheint, stets als Primzahlen gewählt werden.

\emph{Beispiel.} Gegeben sei die Kongruenz $19x\equiv 1 \ (\mathrm{mod.}\ 140)$. Man löse sie zunächst nach dem Modul $2$ und erhält $x\equiv 1\ (\mathrm{mod.}\ 2)$. Setzt man $x=1+2x'$, so wird $38x'\equiv -18\ (\mathrm{mod.}\ 140)$, womit $19x'\equiv -9\ (\mathrm{mod.}\ 70)$ gleichwertig ist. Löst man diese
\SourceLeaf{28}
wiederum nach dem Modul $2$, so ergibt sich $x'\equiv 1\ (\mathrm{mod.}\ 2)$. Mit $x'=1+2x''$ erhält man $38x''\equiv -28\ (\mathrm{mod.}\ 70)$ oder $19x''\equiv -14\ (\mathrm{mod.}\ 35)$. Die Lösung nach $5$ ergibt $x''\equiv 4\ (\mathrm{mod.}\ 5)$, und mit $x''=4+5x'''$ folgt $95x'''\equiv -90\ (\mathrm{mod.}\ 35)$ oder $19x'''\equiv -18\ (\mathrm{mod.}\ 7)$. Daraus folgt schließlich $x'''\equiv 2\ (\mathrm{mod.}\ 7)$, und setzt man $x'''=2+7x''''$, so ergibt sich $x=59+140x''''$. Daher ist $x\equiv 59\ (\mathrm{mod.}\ 140)$ die vollständige Lösung der gegebenen Kongruenz.

\Article{31}

So wie die Wurzel der Gleichung $ax=b$ durch $\frac{b}{a}$ ausgedrückt wird, bezeichnen wir auch jede Wurzel der Kongruenz $ax\equiv b$ durch $\frac{b}{a}$ und fügen zur Unterscheidung den Modul der Kongruenz hinzu. Zum Beispiel bezeichnet $\frac{19}{17}\ (\mathrm{mod.}\ 12)$ jede Zahl, die $\equiv 11\ (\mathrm{mod.}\ 12)$ ist\begingroup\renewcommand{\thefootnote}{*)}\footnote{was entsprechend durch $\frac{11}{1}\ (\mathrm{mod.}\ 12)$ bezeichnet werden kann.}\endgroup. Aus dem Vorhergehenden folgt allgemein, dass $\frac{a}{b}\ (\mathrm{mod.}\ c)$ nichts Wirkliches bezeichnet (oder, wenn man will, etwas Imaginäres), wenn $a,c$ einen gemeinsamen Teiler haben, der $b$ nicht teilt. Außer in diesem Fall hat der Ausdruck $\frac{b}{a}\ (\mathrm{mod.}\ c)$ stets wirkliche Werte, und zwar unendlich viele. Diese sind alle nach $c$ kongruent, wenn $a$ zu $c$ teilerfremd ist, oder nach $\frac{c}{\delta}$, wenn $\delta$ der größte gemeinsame Teiler von $c,a$ ist.

Für diese Ausdrücke gelten fast dieselben Rechenregeln wie für gewöhnliche Brüche. Wir geben hier einige Eigenschaften an, die sich leicht aus dem Vorhergehenden ableiten lassen.

\hangindent=1.7em\hangafter=1
1.\quad Ist nach dem Modul $c$ sowohl $a\equiv\alpha$ als auch $b\equiv\beta$, so sind die Ausdrücke $\frac{a}{b}\ (\mathrm{mod.}\ c)$ und $\frac{\alpha}{\beta}\ (\mathrm{mod.}\ c)$ gleichwertig.

\hangindent=1.7em\hangafter=1
2.\quad $\frac{a\delta}{b\delta}\ (\mathrm{mod.}\ c\delta)$ und $\frac{a}{b}\ (\mathrm{mod.}\ c)$ sind gleichwertig.

\hangindent=1.7em\hangafter=1
3.\quad $\frac{ak}{bk}\ (\mathrm{mod.}\ c)$ und $\frac{a}{b}\ (\mathrm{mod.}\ c)$ sind gleichwertig, wenn $k$ zu $c$ teilerfremd ist.

Viele weitere ähnliche Sätze ließen sich anführen. Da sie aber keine Schwierigkeiten bereiten und für das Folgende nicht besonders nötig sind, gehen wir zu anderem über.

\Needspace{6\baselineskip}
\begin{center}\small\emph{Bestimmung einer Zahl mit gegebenen Resten nach gegebenen Moduln.}\end{center}

\Article{32}

Die Aufgabe, \emph{alle Zahlen zu finden, die nach beliebig vielen gegebenen Moduln gegebene Reste ergeben}, wird im Folgenden sehr nützlich sein und lässt sich aus dem Vorhergehenden leicht lösen. Zunächst seien zwei Moduln $A,B$ gegeben, nach denen die gesuchte Zahl
\SourceLeaf{29}
$z$ zu den Zahlen $a,b$ jeweils kongruent sein soll. Alle Werte von $z$ haben also die Form $Ax+a$, wobei $x$ unbestimmt ist, aber $Ax+a\equiv b\ (\mathrm{mod.}\ B)$ erfüllen muss. Ist nun der größte gemeinsame Teiler von $A,B$ gleich $\delta$, so hat die vollständige Lösung dieser Kongruenz die Form $x\equiv v\ (\mathrm{mod.}\ \frac{B}{\delta})$ oder, gleichbedeutend, $x=v+\frac{kB}{\delta}$, wobei $k$ eine beliebige ganze Zahl bezeichnet. Die Formel $Av+a+\frac{kAB}{\delta}$ umfasst somit alle Werte von $z$, das heißt, $z\equiv Av+a\ (\mathrm{mod.}\ \frac{AB}{\delta})$ ist die vollständige Lösung der Aufgabe. Kommt zu den Moduln $A,B$ ein dritter, $C$, hinzu, nach dem die gesuchte Zahl $z$ $\equiv c$ sein soll, so geht man offenbar genauso vor, da die ersten beiden Bedingungen bereits zu einer einzigen zusammengefasst sind. Ist nämlich der größte gemeinsame Teiler von $\frac{AB}{\delta}.$ und $C$ gleich $=\varepsilon$ und die Lösung der Kongruenz $\frac{AB}{\delta}x+Av+a\equiv c\ (\mathrm{mod.}\ C)$ durch $x\equiv w\ (\mathrm{mod.}\ \frac{C}{\varepsilon})$ gegeben, so ist die Aufgabe durch die Kongruenz $z\equiv\frac{ABw}{\delta}+Av+a\ (\mathrm{mod.}\ \frac{ABC}{\delta\varepsilon})$ vollständig gelöst. Ebenso verfährt man bei beliebig vielen Moduln. Dabei sind $\frac{AB}{\delta}$, $\frac{ABC}{\delta\varepsilon}$ die kleinsten gemeinsamen Vielfachen von $A,B$ beziehungsweise $A,B,C$. Für beliebig viele Moduln $A,B,C$ usw., deren kleinstes gemeinsames Vielfaches $M$ ist, hat daher die vollständige Lösung die Form $z\equiv r\ (\mathrm{mod.}\ M)$. Ist eine der Hilfskongruenzen unlösbar, so muss man schließen, dass die Aufgabe unmöglich ist. Das kann offenbar nicht vorkommen, wenn alle Zahlen $A,B,C$ usw. paarweise teilerfremd sind.

\emph{Beispiel.} Die Zahlen $A,B,C$ und $a,b,c$ seien $504,35,16$ beziehungsweise $17,-4,33$. Dann sind die beiden Bedingungen, dass $z$ $\equiv 17\ (\mathrm{mod.}\ 504)$ und $\equiv -4\ (\mathrm{mod.}\ 35)$ sei, mit der einzigen Bedingung gleichwertig, dass es $\equiv 521\ \mathrm{mod.}\ 2520)$ sei. Zusammen mit $z\equiv 33\ (\mathrm{mod.}\ 16)$ ergibt dies $z\equiv 3041\ (\mathrm{mod.}\ 5040)$.

\Article{33}

Sind die Zahlen $A,B,C$ usw. paarweise teilerfremd, so ist bekanntlich ihr Produkt das kleinste gemeinsame Vielfache aller dieser Zahlen. In diesem Fall sind alle Kongruenzen $z\equiv a\ (\mathrm{mod.}\ A)$, $z\equiv b\ (\mathrm{mod.}\ B)$ usw. offenbar mit der einzigen Kongruenz $z\equiv r\ (\mathrm{mod.}\ R)$ gleichwertig, wobei $R$ das Produkt von $A,B,C$ usw. bezeichnet. Umgekehrt lässt sich daher die einzige Bedingung $z\equiv r\ (\mathrm{mod.}\ R)$ in mehrere zerlegen. Wird $R$ auf irgendeine Weise in paarweise teilerfremde Faktoren $A,B,C$ usw. zerlegt, so erfüllen die Bedingungen $z\equiv r\ (\mathrm{mod.}\ A)$, $z\equiv r\ (\mathrm{mod.}\ B)$, $z\equiv r\ (\mathrm{mod.}\ C)$ usw. die Anforderung vollständig. Diese Bemerkung eröffnet uns einen Weg,
\SourceLeaf{30}
nicht nur eine etwaige Unmöglichkeit der gegebenen Bedingungen sofort zu erkennen, sondern auch die Rechnung bequemer und übersichtlicher durchzuführen.

\Article{34}

Wie oben seien die Bedingungen $z\equiv a\ (\mathrm{mod.}\ A)$, $z\equiv b\ (\mathrm{mod.}\ B)$, $z\equiv c\ (\mathrm{mod.}\ C)$ gegeben. Man zerlege jeden Modul in paarweise teilerfremde Faktoren, $A$ in $A'A''A'''$ usw., $B$ in $B'B''B'''$ usw., und zwar so, dass die Zahlen $A',A''$ usw., $B',B''$ usw. Primzahlen oder Primzahlpotenzen sind. Ist eine der Zahlen $A,B,C$ usw. bereits eine Primzahl oder Primzahlpotenz, so ist für sie keine Zerlegung nötig. Aus dem Vorhergehenden folgt, dass man die gegebenen Bedingungen durch $z\equiv a\ (\mathrm{mod.}\ A')$, $z\equiv a\ (\mathrm{mod.}\ A'')$, $z\equiv a\ (\mathrm{mod.}\ A''')$ usw., $z\equiv b\ (\mathrm{mod.}\ B')$, $z\equiv b\ (\mathrm{mod.}\ B'')$ usw. ersetzen kann. Sind nun nicht alle Zahlen $A,B,C$ usw. paarweise teilerfremd, ist etwa $A$ zu $B$ nicht teilerfremd, so können offenbar die Primteiler von $A,B$ nicht alle verschieden sein. Unter den Faktoren $A',A'',A'''$ usw. muss mindestens einer unter $B',B'',B'''$ usw. entweder einen gleichen Faktor, ein Vielfaches oder einen Teiler haben. Ist \emph{erstens} $A'=B'$, so müssen die Bedingungen $z\equiv a\ (\mathrm{mod.}\ A')$, $z\equiv b\ (\mathrm{mod.}\ B')$ identisch sein, das heißt, es muss $a\equiv b\ (\mathrm{mod.}\ A'\ \text{oder}\ B')$ gelten. Dann kann eine der beiden wegfallen. Gilt dagegen nicht $a\equiv b\ (\mathrm{mod.}\ A')$, so ist die Aufgabe unmöglich. Ist \emph{zweitens} $B'$ ein Vielfaches von $A'$, so muss die Bedingung $z\equiv a\ (\mathrm{mod.}\ A')$ in $z\equiv b\ (\mathrm{mod.}\ B')$ enthalten sein. Das heißt, die aus der letzteren folgende Bedingung $z\equiv b\ (\mathrm{mod.}\ A')$ muss mit der ersteren identisch sein. Also kann die Bedingung $z\equiv a\ (\mathrm{mod.}\ A')$ wegfallen, sofern sie der anderen nicht widerspricht; andernfalls ist die Aufgabe unmöglich. Sind so alle überflüssigen Bedingungen entfernt, so sind offenbar alle verbleibenden Moduln unter $A',A'',A'''$ usw., $B',B'',B'''$ usw. paarweise teilerfremd. Dann ist die Aufgabe sicher möglich, und wir können nach den zuvor gegebenen Regeln vorgehen.

\Article{35}

\emph{Beispiel.} Soll wie oben $z\equiv 17\ (\mathrm{mod.}\ 504)$, $\equiv -4\ (\mathrm{mod.}\ 35)$ und $\equiv 33\ (\mathrm{mod.}\ 16)$ gelten, so lassen sich diese Bedingungen in $z\equiv 17\ (\mathrm{mod.}\ 8)$, $\equiv 17\ (\mathrm{mod.}\ 9)$, $\equiv 17\ (\mathrm{mod.}\ 7)$, $\equiv -4\ (\mathrm{mod.}\ 5)$, $\equiv -4\ (\mathrm{mod.}\ 7)$, $\equiv 33\ (\mathrm{mod.}\ 16)$ zerlegen. Davon können die Bedingungen $z\equiv 17\ (\mathrm{mod.}\ 8)$, $z\equiv 17\ (\mathrm{mod.}\ 7)$ wegfallen, da die erstere in $z\equiv 33\ (\mathrm{mod.}\ 16)$ enthalten und die letztere mit $z\equiv -4\ (\mathrm{mod.}\ 7)$ identisch ist. Es bleiben also

\SourceLeaf{31}\[
z\equiv\left\{\begin{array}{rl}
17 & (\mathrm{mod.}\ 9)\\
-4 & (\mathrm{mod.}\ 5)\\
-4 & (\mathrm{mod.}\ 7)\\
33 & (\mathrm{mod.}\ 16)
\end{array}\right.
\quad\text{daraus folgt}\quad z\equiv 3041\ (\mathrm{mod.}\ 5040).
\]

Meist ist es bequemer, diejenigen verbleibenden Bedingungen, die aus ein und derselben Bedingung hervorgegangen sind, jeweils wieder zusammenzufassen, da dies leicht möglich ist. Sind \emph{zum Beispiel} von den Bedingungen $z\equiv a\ (\mathrm{mod.}\ A')$, $z\equiv a\ (\mathrm{mod.}\ A'')$ usw. einige weggefallen, so erhält man aus den übrigen die Bedingung $z\equiv a$ nach dem Modul, der das Produkt aller verbleibenden Moduln unter $A',A'',A'''$ usw. ist. So ergibt sich in unserem Beispiel aus $z\equiv -4\ (\mathrm{mod.}\ 5)$, $z\equiv -4\ (\mathrm{mod.}\ 7)$ von selbst wieder die ursprüngliche Bedingung $z\equiv -4\ (\mathrm{mod.}\ 35)$. Daraus folgt auch, dass es für die Kürze der Rechnung nicht ganz gleichgültig ist, welche der überflüssigen Bedingungen wegfallen. Diese und andere praktische Verfahren, die man durch Übung viel leichter als durch Regeln lernt, hier darzustellen, ist aber nicht unsere Absicht.

\Article{36}

Sind alle Moduln $A,B,C,D$ usw. paarweise teilerfremd, so ist häufig die folgende Methode vorzuziehen. Man bestimme eine Zahl $\alpha$, die nach $A$ der Eins, nach dem Produkt der übrigen Moduln aber der Null kongruent ist. Oder es sei $\alpha$ ein beliebiger Wert (meist nimmt man am besten den \emph{kleinsten}) des Ausdrucks $\frac{1}{BCD\ \mathrm{usw.}}\ (\mathrm{mod.}\ A)$, multipliziert mit $BCD$ usw. (siehe Artikel 32). Ebenso sei $\beta\equiv 1\ (\mathrm{mod.}\ B)$ und $\equiv 0\ (\mathrm{mod.}\ ACD\ \mathrm{usw.})$, $\gamma\equiv 1\ (\mathrm{mod.}\ C)$ und $\equiv 0\ (\mathrm{mod.}\ ABD\ \mathrm{usw.})$ usw. Wird dann eine Zahl $z$ gesucht, die nach den Moduln $A,B,C,D$ usw. zu den Zahlen $a,b,c,d$ usw. jeweils kongruent ist, so kann man setzen:

\[
z\equiv\alpha a+\beta b+\gamma c+\delta d\ \mathrm{usw.}\ (\mathrm{mod.}\ ABCD\ \mathrm{usw.}).
\]

Denn offenbar ist $\alpha a\equiv a\ (\mathrm{mod.}\ A.)$, während die übrigen Glieder $\beta b,\gamma c$ usw. alle $\equiv 0\ (\mathrm{mod.}\ A)$ sind; also ist $z\equiv a\ (\mathrm{mod.}\ A)$. Der Beweis für die übrigen Moduln verläuft ebenso. Diese Lösung ist der früheren vorzuziehen, wenn mehrere solche Aufgaben zu lösen sind, bei denen die Moduln $A,B,C$ usw. ihre Werte behalten. Dann haben nämlich die Zahlen $\alpha,\beta,\gamma$ usw. konstante Werte. Dies ist bei der Aufgabe der Zeitrechnung der Fall, das Jahr innerhalb der Julianischen Periode zu finden, dessen Indiktion, Goldene Zahl und Sonnenzirkel gegeben sind. Hier ist $A=15$, $B=19$, $C=28$. Da
\SourceLeaf{32}
der Wert des Ausdrucks $\frac{1}{19.28}\ (\mathrm{mod.}\ 15)$ oder $\frac{1}{532}\ (\mathrm{mod.}\ 15)$ gleich $13$ ist, ist $\alpha=6916$. Ebenso findet man für $\beta$ den Wert $4200$ und für $\gamma$ den Wert $4845$. Die gesuchte Zahl ist also der kleinste Rest von $6916a+4200b+4845c$, wobei $a$ die Indiktion, $b$ die Goldene Zahl und $c$ den Sonnenzirkel bezeichnet.

\Needspace{6\baselineskip}
\begin{center}\emph{Lineare Kongruenzen mit mehreren Unbekannten.}\end{center}

\Article{37}

Dies möge über Kongruenzen ersten Grades mit einer einzigen Unbekannten genügen. Es bleiben noch die Kongruenzen zu behandeln, in denen mehrere Unbekannte vorkommen. Eine Darstellung aller Einzelheiten mit voller Strenge würde jedoch viel Raum beanspruchen. Wir wollen hier nicht alles erschöpfend behandeln, sondern nur das anführen, was besonders beachtenswert erscheint. Deshalb beschränken wir die Untersuchung auf wenige Bemerkungen und behalten uns eine ausführlichere Darstellung für eine andere Gelegenheit vor.

1) Wie bei Gleichungen erkennt man auch hier, dass es ebenso viele Kongruenzen wie zu bestimmende Unbekannte geben muss.

2) Es seien also die Kongruenzen

\[
\begin{array}{l@{\qquad}c@{\qquad}r}
ax+by+cz\ldots\equiv f\ (\mathrm{mod.}\ m)&\ldotp\quad\ldotp\quad\ldotp\quad\ldotp\quad\ldotp&(A)\\
a'x+b'y+c'z\ldots\equiv f'&\ldotp\quad\ldotp\quad\ldotp\quad\ldotp\quad\ldotp&(A')\\
a''x+b''y+c''z\ldots\equiv f''&\ldotp\quad\ldotp\quad\ldotp\quad\ldotp\quad\ldotp\quad\ldotp\quad\ldotp&(A'')\\
\mathrm{usw.}&&
\end{array}
\]

gegeben, ebenso viele, wie es Unbekannte $x,y,z$ usw. gibt.

Man bestimme nun Zahlen $\xi,\xi',\xi''$ usw. so, dass

\[\begin{array}{l}
b\xi+b'\xi'+b''\xi''+\mathrm{usw.}=0\\
c\xi+c'\xi'+c''\xi''+\mathrm{usw.}=0\\
\qquad\mathrm{usw.}
\end{array}\]

gilt, und zwar so, dass sie alle ganzzahlig sind und keinen gemeinsamen Faktor haben. Dass dies möglich ist, folgt aus der Theorie der linearen Gleichungen. Ebenso bestimme man $v,v',v''$ usw., $\zeta,\zeta',\zeta''$ usw. so, dass

\[
\begin{array}{l}av+a'v'+a''v''+\mathrm{usw.}=0\\cv+c'v'+c''v''+\mathrm{usw.}=0\\\qquad\mathrm{usw.}\end{array}
\]

\SourceLeaf{33}\[\begin{array}{l}
a\zeta+a'\zeta'+a''\zeta''+\mathrm{usw.}=0\\
b\zeta+b'\zeta'+b''\zeta''+\mathrm{usw.}=0\\
\qquad\mathrm{usw.}\quad\mathrm{usw.}
\end{array}\]

3) Multipliziert man die Kongruenzen $A,A',A''$ usw. mit $\xi,\xi',\xi''$ usw., dann mit $v,v',v'',$ usw., und addiert jeweils die Ergebnisse, so erhält man offenbar die Kongruenzen

\[
\begin{array}{l}(a\xi+a'\xi'+a''\xi''+\mathrm{usw.})x\equiv f\xi+f'\xi'+f''\xi''+\mathrm{usw.}\\(bv+b'v'+b''v''+\mathrm{usw.})y\equiv fv+f'v'+f''v''+\mathrm{usw.}\\(c\zeta+c'\zeta'+c''\zeta''+\mathrm{usw.})z\equiv f\zeta+f'\zeta'+f''\zeta''+\mathrm{usw.}\\\qquad\mathrm{usw.}\end{array}
\]

die wir der Kürze halber so schreiben:

\[
\Sigma(a\xi)x\equiv\Sigma(f\xi),\quad\Sigma(bv)y\equiv\Sigma(fv),\quad\Sigma(c\zeta)z\equiv\Sigma(f\zeta)\ \mathrm{usw.}
\]

4) Nun sind mehrere Fälle zu unterscheiden.

\emph{Erstens}: Sind alle Koeffizienten der Unbekannten, $\Sigma(a\xi),\Sigma(av)$ usw., zum Modul $m$ teilerfremd, so lassen sich diese Kongruenzen nach den zuvor gegebenen Regeln lösen. Die vollständige Lösung der Aufgabe wird durch Kongruenzen der Form $x\equiv p\ (\mathrm{mod.}\ m)$, $y\equiv q\ (\mathrm{mod.}\ m)$ usw. angegeben\begingroup\renewcommand{\thefootnote}{*)}\footnote{Dieser Schluss bedarf eines Beweises, den wir hier jedoch weglassen. Streng genommen folgt aus unserer Untersuchung nur, dass die gegebenen Kongruenzen durch andere Werte der Unbekannten $x,y$ usw. nicht gelöst werden können. Dass diese Werte sie erfüllen, folgt daraus nicht. Es könnte nämlich überhaupt keine Lösung geben. Ein ähnlicher Fehlschluss wird häufig auch bei der Darstellung linearer Gleichungen gemacht.}\endgroup. \emph{Zum Beispiel} seien die Kongruenzen

\[x+3y+z\equiv1,\quad4x+y+5z\equiv7,\quad2x+2y+z\equiv3\ (\mathrm{mod.}\ 8)\]

gegeben. Man findet $\xi=9$, $\xi'=1$, $\xi''=-14$, woraus $-15x\equiv-26$ folgt, also $x\equiv6\ (\mathrm{mod.}\ 8)$. Ebenso erhält man $15y\equiv-4$, $15z\equiv1$ und damit $y\equiv4$, $z\equiv7\ (\mathrm{mod.}\ 8)$.

5) \emph{Zweitens}: Sind nicht alle Koeffizienten $\Sigma(a\xi),\Sigma(bv)$ usw. zum Modul teilerfremd, so seien $\alpha,\beta,\gamma$ usw. die größten gemeinsamen Teiler von $m$ mit $\Sigma(a\xi),\Sigma(bv),\Sigma(c\zeta)$ usw. in der entsprechenden Reihenfolge. Die Aufgabe ist offenbar unmöglich, wenn sie nicht die Zahlen $\Sigma(f\xi),\Sigma(fv),\Sigma(f\zeta)$ usw. jeweils teilen. Gelten diese Bedingungen, so werden die Kongruenzen in (3) durch $x\equiv p\ (\mathrm{mod.}\ \frac{m}{\alpha})$, $y\equiv q\ (\mathrm{mod.}\ \frac{m}{\beta})$, $z\equiv r\ (\mathrm{mod.}\ \frac{m}{\gamma})$ usw. vollständig gelöst. Oder man hat $\alpha$ verschiedene Werte von $x$ (das heißt, nach $m$ inkongruente, nämlich $p,p+\frac{m}{\alpha}\ldots p+\frac{(\alpha-1)m}{\alpha}$),
\SourceLeaf{34}
$\beta$ verschiedene Werte von $y$ usw., die diese Kongruenzen erfüllen. Offenbar finden sich alle Lösungen der gegebenen Kongruenzen, falls es überhaupt welche gibt, unter diesen. Diesen Schluss darf man aber nicht umkehren. Im Allgemeinen erfüllen nicht alle Kombinationen sämtlicher $\alpha$ Werte von $x$ mit allen Werten von $y$ und allen Werten von $z$ usw. die Aufgabe, sondern nur einige. Ihr Zusammenhang lässt sich durch eine oder mehrere Kongruenzen als Bedingungen angeben. Da die vollständige Lösung dieser Aufgabe für das Folgende nicht nötig ist, gehen wir hier nicht näher darauf ein. Es genügt uns, durch ein Beispiel eine Vorstellung davon zu geben.

Gegeben seien die Kongruenzen

\[3x+5y+z\equiv4,\quad2x+3y+2z\equiv7,\quad5x+y+3z\equiv6\ (\mathrm{mod.}\ 12)\]

Hier werden $\xi,\xi',\xi''$, $v,v',v''$, $\zeta,\zeta',\zeta''$ der Reihe nach gleich $=1,-2,1$, $1,1,-1$, $-13,22,-1$. Daraus folgen $4x\equiv-4$, $7y\equiv5$, $28z\equiv96$. Man erhält vier Werte von $x$, nämlich $\equiv2,5,8.\ 11$, einen Wert von $y$, nämlich $\equiv11$, und vier Werte von $z$, nämlich $\equiv0,3,6,9\ (\mathrm{mod.}\ 12)$. Um zu erkennen, welche Kombinationen der Werte von $x$ mit den Werten von $z$ zulässig sind, setzen wir in den gegebenen Kongruenzen für $x,y,z$ der Reihe nach $2+3t$, $11$, $3u$ ein. Sie gehen über in

\[57+9t+3u\equiv0,\quad30+6t+6u\equiv0,\quad15+15t+9u\equiv0\ (\mathrm{mod.}\ 12)\]

Damit sind offenbar die folgenden Kongruenzen gleichwertig:

\[19+3t+u\equiv0,\quad10+2t+2u\equiv0,\quad5+5t+3u\equiv0\ (\mathrm{mod.}\ 4)\]

Die erste verlangt offenbar $u\equiv t+1\ (\mathrm{mod.}\ 4)$. Setzt man diesen Wert in die übrigen Kongruenzen ein, so sind auch diese erfüllt. Folglich müssen die Werte von $x$, nämlich $2,5,8,11$ (die man durch $t\equiv0.\ 1,2,3$ erhält), mit den Werten von $z$, nämlich $z\equiv3,6,9,0$, jeweils in dieser Reihenfolge verbunden werden. Insgesamt gibt es also vier Lösungen:

\[\begin{array}{r@{\;}r@{,\quad}r@{,\quad}r@{,\quad}r@{\ }l}
x\equiv&2&5&8&11&(\mathrm{mod.}\ 12)\\
y\equiv&11&11&11&11&\\
z\equiv&3&6&9&0&
\end{array}\]

\begin{center}$\ast\quad\underset{\ast}{\phantom{\ast}}\quad\ast$\end{center}
